% TOP_PROBLEM: {"id": 2305022, "title": "Research Problems in Function Theory — Problem 5.22", "category_id": 8, "category": {"id": 8, "name": "analysis", "display_name": "Analysis", "description": "Limits, continuity, calculus, and function theory.", "slug": "analysis", "order_index": 8, "created_at": "2026-07-31T15:26:25.671Z"}, "status": "open", "problem_number": "AMR-022-5022", "set_id": 13}
% FIRST_POSTED: 2026-10-02
% ENTRY_KIND: statement_only
% UnsolvedMath ID 2305022; source catalogue code AMR-022-5022
% !TeX program = lualatex
% Definition and sources only; this file is not an LLM solution attempt.
\documentclass[11pt,a4paper]{article}
\usepackage[margin=25mm]{geometry}
\usepackage{fontspec}
\setmainfont{lmroman10-regular.otf}[BoldFont=lmroman10-bold.otf,
 ItalicFont=lmroman10-italic.otf,BoldItalicFont=lmroman10-bolditalic.otf]
\usepackage{amsmath,amssymb,mathtools}
\usepackage{unicode-math}
\setmathfont{latinmodern-math.otf}
\AtBeginDocument{\let\setminus\smallsetminus}
\usepackage{xurl,hyperref}
\hypersetup{colorlinks=true,urlcolor=blue}
\setcounter{secnumdepth}{0}
\setlength{\parindent}{0pt}
\setlength{\parskip}{5pt}
\setlength{\emergencystretch}{3em}
\begin{document}
\section*{Research Problems in Function Theory — Problem 5.22}\label{2305022}
{\small UnsolvedMath ID 2305022 · AMR-022-5022 · Analysis · AMR Open Problem Lists}
\subsection{Definitions and mathematical statement}
Fix $0<p<1$ and let $\mathbb D=\{z\in\mathbb C:|z|<1\}$. The analytic Hardy space $H^p$ consists of holomorphic functions $f$ on $\mathbb D$ for which
\[
\|f\|_{H^p}^p=\sup_{0<r<1}\frac1{2\pi}
\int_0^{2\pi}|f(re^{i\theta})|^p\,d\theta<\infty.
\]
This is a complete quasi-normed space; for $p<1$ the displayed $p$-th-root expression is not generally a norm.
A complex sequence $\lambda=(\lambda_n)_{n\ge0}$ is a coefficient multiplier of $H^p$ into itself if, for every $f(z)=\sum_{n\ge0}a_nz^n\in H^p$, the series
\[
T_\lambda f(z)=\sum_{n\ge0}\lambda_na_nz^n
\]
converges on $\mathbb D$ and defines another member of $H^p$. Characterize exactly all such sequences for each $p\in(0,1)$.
The desired characterization should be an intrinsic condition on $\lambda$, rather than a restatement that the operator works on every input function. Coefficientwise multiplication is intended; this is different from multiplication of functions, which convolves their Taylor coefficients. Neither mere boundedness of $\lambda_n$ nor conditions adequate for multipliers into a different target space may be presumed to answer the self-multiplier problem.
\subsection{Short English statement}
Which sequences can multiply Taylor coefficients of every $H^p$ function, for $0<p<1$, without leaving that same Hardy space?
\subsection{Sources}
\begin{itemize}
\item[S1] ULAM.AI and the UnsolvedMath Contributors, supplied problems.json, record 2305022 (AMR-022-5022), AMR Open Problem Lists. Catalogue identity and source transcription; CC BY 4.0.
\url{https://huggingface.co/datasets/ulamai/UnsolvedMath}.
\item[S2] Hayman - Research Problems in Function Theory (2018), Problem 5.22. Original source indexed by AMR-022-5022.
\url{https://arxiv.org/abs/1809.07200}.
\end{itemize}
\end{document}
